% lecture11.tex - 特征值计算与特征基
% 第 11 章: 特征值计算与特征基

\chapter{特征值计算与特征基}
\label{ch:eigenbasis}

\begin{abstract}
在本章中, 我们将学习计算 $2 \times 2$ 矩阵特征值的快捷方法, 并熟悉特征基 (eigenbasis) 与对角化的概念。
\end{abstract}

% ============================================
\section{$2 \times 2$ 矩阵的快捷算法}
% ============================================

\begin{theorem}[快捷公式]
对于矩阵 $\mA = \twomat{a}{b}{c}{d}$:

\textbf{特征值之和:}
\[
\lambda_1 + \lambda_2 = a + d = \tr(\mA)
\]

\textbf{特征值之积:}
\[
\lambda_1 \cdot \lambda_2 = ad - bc = \det(\mA)
\]
\end{theorem}

\begin{intuition}
已知 $m = \lambda_1 + \lambda_2$ 与 $p = \lambda_1 \lambda_2$, 可求出 $\lambda_{1,2}$:
\[
\lambda = \frac{m}{2} \pm \sqrt{\left(\frac{m}{2}\right)^2 - p}
\]
或者: 找出和为 $m$、积为 $p$ 的两个数。
\end{intuition}

\begin{example}
$\mA = \twomat{3}{1}{4}{1}$

迹: $m = 3 + 1 = 4$

行列式: $p = 3 - 4 = -1$

特征值: 满足 $x + y = 4$ 且 $xy = -1$ 的两个数:
\[
\lambda = 2 \pm \sqrt{4 - (-1)} = 2 \pm \sqrt{5}
\]
\end{example}

% ============================================
\section{特征基}
% ============================================

\begin{definition}[特征基]
若矩阵的特征向量可以构成一组\vocab{基 (basis)} (即有 $n$ 个线性无关 (linear independence) 的特征向量), 则该基称为\vocab{特征基 (eigenbasis)}。
\end{definition}

\begin{theorem}
在特征基下, 变换矩阵变为\vocab{对角的 (diagonal)}:
\[
[\mA]_{\text{eigenbasis}} = \mD = \threemat{\lambda_1}{0}{0}{0}{\lambda_2}{0}{0}{0}{\ddots}
\]
\end{theorem}

\begin{intuition}
为什么是对角的?

在特征基下, 每个基向量只被缩放:
\begin{align*}
\mA \vv_1 &= \lambda_1 \vv_1 \to \text{第一列: } \threevec{\lambda_1}{0}{0} \\
\mA \vv_2 &= \lambda_2 \vv_2 \to \text{第二列: } \threevec{0}{\lambda_2}{0}
\end{align*}
\end{intuition}

% ============================================
\section{对角化}
% ============================================

\begin{definition}[对角化]
若矩阵 $\mA$ 可写为:
\[
\mA = \mP \mD \mP\inv
\]
其中 $\mD$ 为对角矩阵 (diagonal matrix), $\mP$ 为以特征向量为列的矩阵, 则称 $\mA$ \vocab{可对角化 (diagonalizable)}。
\end{definition}

\begin{theorem}[可对角化条件]
$n \times n$ 矩阵可对角化当且仅当它有 $n$ 个线性无关的特征向量。
\end{theorem}

\begin{example}
$\mA = \twomat{3}{1}{0}{2}$, 特征值为 $\lambda_1 = 3$, $\lambda_2 = 2$

特征向量: $\vv_1 = \twovec{1}{0}$, $\vv_2 = \twovec{-1}{1}$

\[
\mP = \twomat{1}{-1}{0}{1}, \quad \mD = \twomat{3}{0}{0}{2}
\]

验证: $\mP\mD\mP\inv = \mA$ \checkmark
\end{example}

% ============================================
\section{用对角化计算矩阵幂}
% ============================================

\begin{theorem}
若 $\mA = \mP\mD\mP\inv$:
\[
\mA^n = \mP \mD^n \mP\inv
\]
而 $\mD^n$ 非常简单:
\[
\mD^n = \threemat{\lambda_1^n}{0}{0}{0}{\lambda_2^n}{0}{0}{0}{\ddots}
\]
\end{theorem}

\begin{example}
计算上一例中 $\mA^{100}$:
\[
\mA^{100} = \mP \twomat{3^{100}}{0}{0}{2^{100}} \mP\inv
\]
若不用对角化, 需要做 100 次矩阵乘法!
\end{example}

\begin{practical}
\textbf{应用: 斐波那契数列}

斐波那契数列: $F_{n+1} = F_n + F_{n-1}$

矩阵形式:
\[
\twomat{1}{1}{1}{0}^n = \twomat{F_{n+1}}{F_n}{F_n}{F_{n-1}}
\]

利用对角化, 可得到闭式公式:
\[
F_n = \frac{1}{\sqrt{5}}\left[\left(\frac{1+\sqrt{5}}{2}\right)^n - \left(\frac{1-\sqrt{5}}{2}\right)^n\right]
\]
\end{practical}

% ============================================
\section{不可对角化的矩阵}
% ============================================

\begin{example}
错切矩阵 $\mA = \twomat{1}{1}{0}{1}$:

特征多项式: $(1-\lambda)^2 = 0$

特征值: $\lambda = 1$ (重根)

特征向量: 仅 $\twovec{1}{0}$ (一维)

该矩阵不可对角化!
\end{example}

\begin{warning}
重特征值本身不一定是问题。当独立特征向量的个数少于该特征值的重数时, 问题才会出现。
\end{warning}

% ============================================
\section{对称矩阵}
% ============================================

\begin{theorem}[谱定理]
对称矩阵 ($\mA = \mA\trans$):
\begin{enumerate}
    \item 有实特征值
    \item 属于不同特征值的特征向量相互正交
    \item 总可对角化 (且有一组正交基)
\end{enumerate}
\end{theorem}

\begin{intuition}
对称矩阵是``表现良好''的矩阵。它们总能化为对角形, 且特征向量相互垂直——就像椭圆的主轴一样。
\end{intuition}

% ============================================
\section{练习}
% ============================================

\begin{exercise}
用快捷公式求下列矩阵的特征值:
\[
\mA = \twomat{5}{2}{2}{2}
\]
\end{exercise}

\begin{exercise}
将矩阵 $\mA = \twomat{2}{1}{1}{2}$ 对角化。
\end{exercise}

\begin{exercise}
利用对角化计算 $\mA^{10}$:
\[
\mA = \twomat{1}{1}{0}{2}
\]
\end{exercise}

\begin{exercise}
下列矩阵是否可对角化?
\[
\mA = \twomat{2}{1}{0}{2}
\]
\end{exercise}

\begin{exercise}[挑战]
证明 $\mA$ 与 $\mA\trans$ 有相同的特征值。
\end{exercise}

\begin{problem}
兔子与狐狸的数量由下式建模:
\[
\twovec{R_{n+1}}{F_{n+1}} = \twomat{1.1}{-0.4}{0.2}{0.8}\twovec{R_n}{F_n}
\]
分析种群的长期行为。
\end{problem}
